\section{Mở đầu}

\subsection{Bối cảnh nghiên cứu}

Học sâu trong vài năm gần đây không còn giới hạn ở dữ liệu ảnh mà đã mở rộng
mạnh sang xử lý âm thanh. Khác với ảnh -- vốn là một ma trận hai chiều tương
đối ổn định -- tín hiệu âm thanh là một chuỗi biến thiên theo thời gian, độ dài
không cố định và dễ bị nhiễu môi trường chi phối. Chính những đặc điểm này khiến
việc nhận diện âm thanh tự động trở thành một bài toán vừa khó vừa đáng làm.

Trong không gian đô thị, âm thanh hiện diện dày đặc: tiếng máy điều hòa, còi xe,
tiếng trẻ chơi đùa, chó sủa, máy khoan, động cơ chạy không tải, tiếng súng, búa
khoan, còi báo động và nhạc đường phố. Mỗi loại đều mang một phần thông tin về
môi trường xung quanh. Một hệ thống biết phân loại chúng có thể phục vụ giám sát
môi trường, hỗ trợ thành phố thông minh, cảnh báo sự kiện bất thường hoặc làm
đầu vào cho các hệ thống phân tích ngữ cảnh ở mức cao hơn.

Bài toán \textit{urban sound classification} vì vậy có giá trị cả về ứng dụng
lẫn học thuật. Hướng tiếp cận phổ biến trong học sâu là không làm việc trực tiếp
với dạng sóng thô mà chuyển tín hiệu sang miền thời gian -- tần số thông qua
spectrogram hoặc Mel spectrogram, rồi đưa biểu diễn hai chiều này vào mạng nơ-ron
tích chập. Đề tài này đi theo đúng hướng đó.

\subsection{Động lực thực hiện đề tài}

Đề tài được lựa chọn trước hết vì nó là một bài toán học sâu nằm \textit{ngoài}
miền ảnh quen thuộc. Làm việc với âm thanh buộc người thực hiện phải hiểu rõ vai
trò của tiền xử lý và biểu diễn đặc trưng -- những bước mà ở bài toán ảnh thường
được xem nhẹ. Đây là cơ hội tốt để nhìn lại toàn bộ một pipeline học sâu thay vì
chỉ tập trung vào kiến trúc mô hình.

Bộ dữ liệu UrbanSound8K phù hợp với phạm vi một bài tập lớn cá nhân: đủ lớn để
kết quả có ý nghĩa nhưng vẫn vừa sức huấn luyện trên tài nguyên hạn chế, lại được
cộng đồng nghiên cứu sử dụng rộng rãi nên có nhiều mốc tham chiếu. Cuối cùng, đề
tài cho phép triển khai trọn vẹn một quy trình -- từ phân tích dữ liệu, tiền xử
lý, xây dựng \texttt{Dataset}/\texttt{DataLoader}, huấn luyện, đánh giá đến viết
báo cáo -- nên rất thích hợp để tổng hợp kiến thức của môn học.

\subsection{Mục tiêu nghiên cứu}

Mục tiêu tổng quát là xây dựng và đánh giá một hệ thống phân loại âm thanh đô thị
trên UrbanSound8K, sử dụng biểu diễn log-Mel spectrogram và các mô hình học sâu.
Cụ thể, đề tài hướng tới các nội dung sau:

\begin{itemize}[leftmargin=1.4em,itemsep=2pt]
  \item Phân tích đặc điểm của UrbanSound8K, đặc biệt là các yếu tố không đồng
        nhất trong dữ liệu gốc, để rút ra quyết định tiền xử lý.
  \item Thiết kế một pipeline tiền xử lý thống nhất, đưa toàn bộ dữ liệu về cùng
        một chuẩn đầu vào và trích xuất log-Mel spectrogram.
  \item Triển khai một mô hình CNN làm mốc cơ sở, một mô hình ResNet18 được điều
        chỉnh cho đầu vào spectrogram một kênh, và khảo sát tác động của
        augmentation trong miền spectrogram.
  \item Kết hợp các mô hình đã huấn luyện theo hướng ensemble nhằm tận dụng tính
        bù trừ lỗi giữa chúng.
  \item Đánh giá và phân tích kết quả bằng Accuracy, Macro-F1, classification
        report và confusion matrix, từ đó so sánh công bằng giữa các cấu hình.
\end{itemize}

Cần nói rõ ngay từ đầu rằng đề tài không đặt mục tiêu vượt các kết quả tốt nhất
đã công bố trên UrbanSound8K. Trong khuôn khổ một bài tập lớn, câu hỏi thực tế
hơn là: với cùng một pipeline rõ ràng và tài nguyên vừa phải, các lựa chọn về
kiến trúc và augmentation ảnh hưởng đến kết quả ra sao, và ta rút ra được điều gì
từ những khác biệt đó.

\subsection{Cấu trúc báo cáo}

Phần còn lại của báo cáo được tổ chức như sau. Mục~\ref{sec:problem} mô tả bài
toán và bộ dữ liệu, kèm kết quả phân tích dữ liệu khám phá (EDA).
Mục~\ref{sec:litreview} điểm lại các hướng tiếp cận cho phân loại âm thanh.
Mục~\ref{sec:method} trình bày phương pháp: pipeline tiền xử lý, các mô hình và
chiến lược huấn luyện. Mục~\ref{sec:expdesign} mô tả thiết kế thực nghiệm.
Mục~\ref{sec:results} trình bày, so sánh và phân tích kết quả của bốn cấu hình.
Mục~\ref{sec:conclusion} tổng kết và nêu hướng phát triển.
